% GENERATED FILE. Edit canonical lessons, then run npm run generate:books.
% canonical-source-hash: fnv1a64:49688ab5324e6ce8
% canonical-lessons: JA-C361-moukeru, JA-C361-moukaru, JA-C361-rieki, JA-C361-son, JA-C361-tsubureru, JA-R361-review-357-361

\chapter{Profit and Loss}
\label{ch:ja-profit-and-loss}
\hlchaptermodality{361}

\begin{chapteropening}
\textbf{By the end of this chapter:} \emph{I can say whether a shop would make or lose money.}

Complete the last lesson of chapter 361: I can say whether a shop would make or lose money.
\end{chapteropening}

\section[\texorpdfstring{\ja{もうける} (mōkeru)}{もうける (mōkeru)}]{\ja{もうける} (mōkeru) — to make a profit}

\label{lesson:JA-C361-moukeru}



\begin{quote}
Before the new one: say the Japanese for to be saved up, then the Japanese for to keep down.
\end{quote}

\subsection*{You'll want to know: \ja{もうける}}
\textbf{\ja{もうける}} — \emph{mōkeru} — \textquotedblleft{}to make a profit\textquotedblright{}.

\begin{quote}
\textbf{\ja{よく} \ja{うれたら、みせは} \ja{もうけるでしょう。}} — \emph{yoku uretara, mise wa mōkeru deshō} — \textquotedblleft{}If it sells well, the shop will make a profit.\textquotedblright{}
\end{quote}

Polite \textbf{\ja{もうけます}}, \emph{mōkemasu}; an ichidan verb. Past (polite): \textbf{\ja{みせが} \ja{もうけました。}}, \emph{mise ga mōkemashita}.

\begin{cousinweb}
\textbf{\ja{もうける}} (\emph{mōkeru}) is a native Japanese word (wago), inherited rather than borrowed from Chinese, and its older origin is not settled. Its kanji was attached to it later as a kun'yomi, a native reading given to a Chinese character (JMdict, KANJIDIC2).
\end{cousinweb}

\begin{grammarlens}[title={What you've built}]
One more word for saying what happens if.
\end{grammarlens}

\subsection*{Your turn}
\begin{itemize}
\raggedright
  \item \emph{Say it:} \emph{mōkeru}
  \item \emph{Say it:} \emph{mōkeru}, once more
  \item \emph{Say it:} \emph{yoku uretara, mise wa mōkeru deshō} — If it sells well, the shop will make a profit.
  \item \emph{Recall:} say the Japanese for to be saved up, then the Japanese for to keep down, then say \emph{mōkeru} again
\end{itemize}

\begin{tcolorbox}[breakable,colback=teal!4,colframe=teal!35!black,title={Before you move on}]
What does \ja{もうける} mean? (\textquotedblleft{}To make a profit\textquotedblright{}.) Say it once more.
\end{tcolorbox}
\section[\texorpdfstring{\ja{もうかる} (mōkaru)}{もうかる (mōkaru)}]{\ja{もうかる} (mōkaru) — to be profitable; to pay}

\label{lesson:JA-C361-moukaru}



\begin{quote}
Before the new one: say the Japanese for to keep down, then the Japanese for to make a profit.
\end{quote}

\subsection*{You'll want to know: \ja{もうかる}}
\textbf{\ja{もうかる}} — \emph{mōkaru} — \textquotedblleft{}to be profitable; to pay\textquotedblright{}.

\begin{quote}
\textbf{\ja{みせが} \ja{もうかったら、てんいんを} \ja{ふやします。}} — \emph{mise ga mōkattara, ten'in o fuyashimasu} — \textquotedblleft{}If the shop does well, we will take on more people.\textquotedblright{}
\end{quote}

Polite \textbf{\ja{もうかります}}, \emph{mōkarimasu}; a godan verb. Past (polite): \textbf{\ja{みせが} \ja{もうかりました。}}, \emph{mise ga mōkarimashita}.

\begin{cousinweb}
\textbf{\ja{もうかる}} (\emph{mōkaru}) is a native Japanese word (wago), inherited rather than borrowed from Chinese, and its older origin is not settled. Its kanji was attached to it later as a kun'yomi, a native reading given to a Chinese character (JMdict, KANJIDIC2).
\end{cousinweb}

\begin{grammarlens}[title={What you've built}]
One more word for saying what happens if.
\end{grammarlens}

\subsection*{Your turn}
\begin{itemize}
\raggedright
  \item \emph{Say it:} \emph{mōkaru}
  \item \emph{Say it:} \emph{mōkaru}, once more
  \item \emph{Say it:} \emph{mise ga mōkattara, ten'in o fuyashimasu} — If the shop does well, we will take on more people.
  \item \emph{Recall:} say the Japanese for to keep down, then the Japanese for to make a profit, then say \emph{mōkaru} again
\end{itemize}

\begin{tcolorbox}[breakable,colback=teal!4,colframe=teal!35!black,title={Before you move on}]
What does \ja{もうかる} mean? (\textquotedblleft{}To be profitable; to pay\textquotedblright{}.) Say it once more.
\end{tcolorbox}
\section[\texorpdfstring{\ja{りえき} (rieki)}{りえき (rieki)}]{\ja{りえき} (rieki) — a profit; gains}

\label{lesson:JA-C361-rieki}



\begin{quote}
Before the new one: say the Japanese for to make a profit, then the Japanese for to be profitable.
\end{quote}

\subsection*{You'll want to know: \ja{りえき}}
\textbf{\ja{りえき}} — \emph{rieki} — \textquotedblleft{}a profit; gains\textquotedblright{}.

\begin{quote}
\textbf{\ja{りえきが} \ja{あったら、みんなで} \ja{わけます。}} — \emph{rieki ga attara, minna de wakemasu} — \textquotedblleft{}If there is a profit, we share it among us.\textquotedblright{}
\end{quote}

\begin{cousinweb}
\textbf{\ja{りえき}} (\emph{rieki}) is Sino-Japanese (kango): two Chinese characters read with their on'yomi, readings borrowed from Chinese. They are \emph{ri}, \textquotedblleft{}advantage\textquotedblright{} (Middle Chinese \emph{lijH}) and \emph{eki}, \textquotedblleft{}benefit\textquotedblright{} (\emph{'jiek}) (KANJIDIC2; Middle Chinese in Baxter's transcription).
\end{cousinweb}

\begin{grammarlens}[title={What you've built}]
One more word for saying what happens if.
\end{grammarlens}

\subsection*{Your turn}
\begin{itemize}
\raggedright
  \item \emph{Say it:} \emph{rieki}
  \item \emph{Say it:} \emph{rieki}, once more
  \item \emph{Say it:} \emph{rieki ga attara, minna de wakemasu} — If there is a profit, we share it among us.
  \item \emph{Recall:} say the Japanese for to make a profit, then the Japanese for to be profitable, then say \emph{rieki} again
\end{itemize}

\begin{tcolorbox}[breakable,colback=teal!4,colframe=teal!35!black,title={Before you move on}]
What does \ja{りえき} mean? (\textquotedblleft{}A profit; gains\textquotedblright{}.) Say it once more.
\end{tcolorbox}
\section[\texorpdfstring{\ja{そん} (son)}{そん (son)}]{\ja{そん} (son) — a loss; a disadvantage}

\label{lesson:JA-C361-son}



\begin{quote}
Before the new one: say the Japanese for to be profitable, then the Japanese for a profit.
\end{quote}

\subsection*{You'll want to know: \ja{そん}}
\textbf{\ja{そん}} — \emph{son} — \textquotedblleft{}a loss; a disadvantage\textquotedblright{}.

\begin{quote}
\textbf{\ja{いま} \ja{うったら、そんを} \ja{するでしょう。}} — \emph{ima uttara, son o suru deshō} — \textquotedblleft{}If you sell now, you will probably make a loss.\textquotedblright{}
\end{quote}

\begin{cousinweb}
\textbf{\ja{そん}} (\emph{son}) is Sino-Japanese (kango): one Chinese character read with its on'yomi, reading borrowed from Chinese. It is \emph{son}, \textquotedblleft{}loss\textquotedblright{} (Middle Chinese \emph{swonX}) (KANJIDIC2; Middle Chinese in Baxter's transcription).
\end{cousinweb}

\begin{grammarlens}[title={What you've built}]
One more word for saying what happens if.
\end{grammarlens}

\subsection*{Your turn}
\begin{itemize}
\raggedright
  \item \emph{Say it:} \emph{son}
  \item \emph{Say it:} \emph{son}, once more
  \item \emph{Say it:} \emph{ima uttara, son o suru deshō} — If you sell now, you will probably make a loss.
  \item \emph{Recall:} say the Japanese for to be profitable, then the Japanese for a profit, then say \emph{son} again
\end{itemize}

\begin{tcolorbox}[breakable,colback=teal!4,colframe=teal!35!black,title={Before you move on}]
What does \ja{そん} mean? (\textquotedblleft{}A loss; a disadvantage\textquotedblright{}.) Say it once more.
\end{tcolorbox}
\section[\texorpdfstring{\ja{つぶれる} (tsubureru)}{つぶれる (tsubureru)}]{\ja{つぶれる} (tsubureru) — to go out of business; to be crushed}

\label{lesson:JA-C361-tsubureru}



\begin{quote}
Before the new one: say the Japanese for a profit, then the Japanese for a loss.
\end{quote}

\subsection*{You'll want to know: \ja{つぶれる}}
\textbf{\ja{つぶれる}} — \emph{tsubureru} — \textquotedblleft{}to go out of business; to be crushed\textquotedblright{}.

\begin{quote}
\textbf{\ja{きゃくが} \ja{すくなく} \ja{なったら、みせは} \ja{つぶれるでしょう。}} — \emph{kyaku ga sukunaku nattara, mise wa tsubureru deshō} — \textquotedblleft{}If the customers fall away, the shop will probably go under.\textquotedblright{}
\end{quote}

Polite \textbf{\ja{つぶれます}}, \emph{tsuburemasu}; an ichidan verb. Past (polite): \textbf{\ja{みせが} \ja{つぶれました。}}, \emph{mise ga tsuburemashita}.

\begin{cousinweb}
\textbf{\ja{つぶれる}} (\emph{tsubureru}) is a native Japanese word (wago), inherited rather than borrowed from Chinese, and its older origin is not settled. Its kanji was attached to it later as a kun'yomi, a native reading given to a Chinese character (JMdict, KANJIDIC2).
\end{cousinweb}

\begin{grammarlens}[title={What you've built}]
One more word for saying what happens if.
\end{grammarlens}

\subsection*{Your turn}
\begin{itemize}
\raggedright
  \item \emph{Say it:} \emph{tsubureru}
  \item \emph{Say it:} \emph{tsubureru}, once more
  \item \emph{Say it:} \emph{kyaku ga sukunaku nattara, mise wa tsubureru deshō} — If the customers fall away, the shop will probably go under.
  \item \emph{Recall:} say the Japanese for a profit, then the Japanese for a loss, then say \emph{tsubureru} again
\end{itemize}

\begin{tcolorbox}[breakable,colback=teal!4,colframe=teal!35!black,title={Before you move on}]
What does \ja{つぶれる} mean? (\textquotedblleft{}To go out of business; to be crushed\textquotedblright{}.) Say it once more.
\end{tcolorbox}
\section[(dialogue)]{Looking back: from Contracts to Profit and Loss}

\label{lesson:JA-R361-review-357-361}



\begin{quote}
Say the Japanese for to sign a contract and a warranty.
\end{quote}

\subsection*{Your turn}
Say these aloud:

\begin{itemize}
\raggedright
  \item to sign a contract, a warranty, the contents, leeway, to decide officially
  \item rent, a landlord, a bank account, to pay into an account, a bank balance
  \item a debt, interest on money, to pay for someone, a refund, a sum of money
  \item income, expenditure, living costs, to be saved up, to keep down
  \item to make a profit, to be profitable, a profit, a loss, to go out of business
\end{itemize}

\begin{tcolorbox}[breakable,colback=teal!4,colframe=teal!35!black,title={Before you move on}]
To sign a contract? (\textbf{\ja{けいやくする}}, \emph{keiyaku suru}.) A warranty? (\textbf{\ja{ほしょう}}, \emph{hoshō}.) The contents? (\textbf{\ja{ないよう}}, \emph{naiyō}.) Leeway? (\textbf{\ja{ゆとり}}, \emph{yutori}.) To decide officially? (\textbf{\ja{けっていする}}, \emph{kettei suru}.) Rent? (\textbf{\ja{やちん}}, \emph{yachin}.) A landlord? (\textbf{\ja{おおや}}, \emph{ōya}.) A bank account? (\textbf{\ja{こうざ}}, \emph{kōza}.) To pay into an account? (\textbf{\ja{ふりこむ}}, \emph{furikomu}.) A bank balance? (\textbf{\ja{ざんだか}}, \emph{zandaka}.) A debt? (\textbf{\ja{しゃっきん}}, \emph{shakkin}.) Interest on money? (\textbf{\ja{りし}}, \emph{rishi}.) To pay for someone? (\textbf{\ja{たてかえる}}, \emph{tatekaeru}.) A refund? (\textbf{\ja{へんきん}}, \emph{henkin}.) A sum of money? (\textbf{\ja{きんがく}}, \emph{kingaku}.) Income? (\textbf{\ja{しゅうにゅう}}, \emph{shūnyū}.) Expenditure? (\textbf{\ja{ししゅつ}}, \emph{shishutsu}.) Living costs? (\textbf{\ja{せいかつひ}}, \emph{seikatsuhi}.) To be saved up? (\textbf{\ja{たまる}}, \emph{tamaru}.) To keep down? (\textbf{\ja{おさえる}}, \emph{osaeru}.) To make a profit? (\textbf{\ja{もうける}}, \emph{mōkeru}.) To be profitable? (\textbf{\ja{もうかる}}, \emph{mōkaru}.) A profit? (\textbf{\ja{りえき}}, \emph{rieki}.) A loss? (\textbf{\ja{そん}}, \emph{son}.) To go out of business? (\textbf{\ja{つぶれる}}, \emph{tsubureru}.)
\end{tcolorbox}
